\section{Introduction}

Modern programmable mixed-signal platforms have become an attractive foundation for low-cost scientific instrumentation because they integrate configurable analog resources together with analog-to-digital and digital-to-analog converters, programmable digital filters, direct memory access, and embedded processors within a single device. However, the configurable analog building blocks available in these platforms generally provide lower precision, larger offsets, and reduced matching than dedicated analog integrated circuits, limiting the achievable performance of the analog front end.

Rather than addressing these limitations through higher-grade analog components, manual trimming, or additional external circuitry, the proposed work follows the principles of digitally assisted analog design introduced by Murmann and Boser \cite{Murmann2004} and further developed by Murmann \cite{Murmann2006}. In this work, a \emph{Digitally Assisted Analog Front-End} (DA-AFE) denotes a system in which embedded digital resources non-intrusively assist the analog circuitry to improve its effective operating conditions while leaving the acquisition signal path unchanged during normal operation. Although digitally assisted analog design has primarily been investigated at the integrated-circuit level, the present work extends the same philosophy to low-cost scientific instrumentation through a geophone acquisition platform.

Low-cost multichannel geophone acquisition systems provide a representative application for this DA-AFE. Surface-wave methods such as Multichannel Analysis of Surface Waves (MASW) estimate Rayleigh-wave dispersion to infer shear-wave velocity profiles for geotechnical site characterization \cite{Park1999,Foti2014}. Low-cost embedded geophone interfaces have already been demonstrated \cite{Kafadar2020,Wijayaraja2024}, and dedicated low-frequency conditioning circuits continue to be reported \cite{Deng2025}, but in multistage analog conditioning chains, accumulated amplifier and reference offsets progressively reduce the available dynamic range and, at high gain, may saturate intermediate stages before the low-amplitude seismic signal reaches the acquisition converter. Improving the effective operating conditions of the analog front end is therefore a precondition for exploiting the available acquisition range.

The proposed DA-AFE forms part of a low-cost reconfigurable Internet-of-Things (IoT) geophone acquisition platform \cite{Guevara2015} based on SM-24 moving-coil geophones \cite{SM24}. The PSoC 5LP integrates programmable analog blocks, voltage DACs (VDACs), a delta-sigma ADC, programmable digital filters, DMA resources, and an Arm Cortex-M3 processor \cite{InfineonPSoC}, providing all resources required to implement the proposed DA-AFE. As a proof of feasibility, this work implements an automatic foreground DC-offset calibration that sequentially recenters the programmable-gain amplifier (PGA), band-pass filter (BP), summing stage (SUM), and low-pass filter (LP) before acquisition.

Previous offset-compensation techniques generally follow three approaches: numerical estimation and digital correction \cite{InfineonAN68403,Lopin2018}, physical trimming of individual analog stages \cite{Mohan2020}, or autozero and chopper stabilization techniques \cite{EnzTemes1996, Liu2020}. In contrast, the proposed DA-AFE employs embedded digital resources only during foreground calibration, after which the analog signal chain operates conventionally throughout data acquisition.

The contributions of this work are threefold: (i) the proposal of a Digitally Assisted Analog Front-End (DA-AFE) that exploits embedded mixed-signal resources to improve the effective operating conditions of configurable analog hardware while preserving hardware simplicity; (ii) its demonstration through automatic foreground DC-offset calibration of a complete low-cost geophone conditioning chain; and (iii) an experimental evaluation on two independently assembled acquisition boards under multiple gain configurations using a single controller configuration across all reported conditions.

The remainder of this paper is organized as follows. Section II presents the proposed DA-AFE. Section III describes the foreground calibration method. Section IV reports the experimental validation, followed by the discussion and conclusion.
